\documentclass[11pt,a4paper]{article}

\usepackage[utf8]{inputenc}
\usepackage[T1]{fontenc}
\usepackage{lmodern}
\usepackage[margin=0.7in]{geometry}
\usepackage[table]{xcolor}
\usepackage{array}
\usepackage{tabularx}
\usepackage{booktabs}
\usepackage{enumitem}
\usepackage{titlesec}
\usepackage{fancyhdr}
\usepackage{lastpage}
\usepackage{hyperref}

\definecolor{TextInk}{HTML}{233142}
\definecolor{HeadingBlue}{HTML}{1D3557}
\definecolor{AccentTeal}{HTML}{2A7F92}
\definecolor{AccentGold}{HTML}{C58B4E}
\definecolor{RuleGray}{HTML}{D6E0E5}
\definecolor{TableStripe}{HTML}{F3F8FA}
\definecolor{SoftTint}{HTML}{F7F4EE}
\definecolor{SoftBlue}{HTML}{EEF6F8}

\hypersetup{
  colorlinks=true,
  linkcolor=HeadingBlue,
  urlcolor=AccentTeal,
  pdftitle={IB Geography SL Flashcards and Summary - 2026-04-24},
  pdfauthor={IB Geography SL Preparation}
}

\color{TextInk}
\setlength{\parskip}{0.45em}
\setlength{\parindent}{0pt}
\renewcommand{\arraystretch}{1.25}
\setlist[itemize]{leftmargin=1.3em,itemsep=2pt,topsep=3pt}

\newcolumntype{Y}{>{\raggedright\arraybackslash}X}
\newcolumntype{L}[1]{>{\raggedright\arraybackslash}p{#1}}

\titleformat{\section}
  {\normalfont\Large\sffamily\bfseries\color{HeadingBlue}}
  {}{0pt}{}
  [\vspace{0.1em}{\color{AccentGold}\titlerule[0.8pt]}]

\titleformat{\subsection}
  {\normalfont\large\sffamily\bfseries\color{AccentTeal}}
  {}{0pt}{}

\pagestyle{fancy}
\setlength{\headheight}{14pt}
\fancyhf{}
\fancyhead[L]{\sffamily\color{AccentTeal}April 24 Workbook}
\fancyhead[R]{\sffamily\color{HeadingBlue}IB Geography SL}
\fancyfoot[C]{\sffamily\color{HeadingBlue}\thepage\ of \pageref{LastPage}}
\renewcommand{\headrulewidth}{0.4pt}

\newcommand{\card}[2]{%
\noindent\fcolorbox{RuleGray}{#1}{%
  \begin{minipage}[t][0.155\textheight][t]{0.455\textwidth}
  \sffamily\textbf{\textcolor{HeadingBlue}{#2}}\par
  \vspace{0.3em}
  Front term or question:\par
  \vspace{1.15cm}
  Back answer: definition + cause/process + consequence + example hook\par
  \vspace{1.15cm}
  \end{minipage}
}%
}

\begin{document}

\begin{center}
{\sffamily\bfseries\color{HeadingBlue}\LARGE April 24 Flashcards And Summary Workbook}\\[0.25em]
{\sffamily\color{AccentTeal}\large Urban Environments + Population}
\end{center}

\noindent\colorbox{SoftBlue}{%
\begin{minipage}{0.965\textwidth}
\sffamily
\textbf{\textcolor{HeadingBlue}{Use after the April 24 study pack.}}
The plan requires ten flashcards or a one-page summary from memory. This
workbook gives both formats. Complete one format first; use the other only if
there is extra time.
\end{minipage}
}

\section{Ten Flashcards}

\subsection{Recommended Card Fronts}

Choose ten cards total. A balanced set is five Urban environments cards and
five Population cards. If you choose only five Population cards, include both
\textbf{dependency ratio} and \textbf{demographic dividend}.

\begin{tabularx}{\textwidth}{L{0.08\textwidth}Y Y}
\toprule
\textbf{No.} & \textbf{Urban environments front bank} & \textbf{Population front bank} \\
\midrule
1 & Urbanization & Population distribution \\
2 & Suburbanization & Natural increase \\
3 & Gentrification & Migration \\
4 & Deindustrialization & Population pyramid \\
5 & Urban environmental stress & Dependency ratio \\
6 & Sustainable urban management & Demographic dividend \\
7 & Urban social stress & Population policy \\
\bottomrule
\end{tabularx}

\subsection{Printable Card Grid}

\card{SoftTint}{Card 1}
\hfill
\card{SoftBlue}{Card 2}

\vspace{0.55em}

\card{SoftBlue}{Card 3}
\hfill
\card{SoftTint}{Card 4}

\vspace{0.55em}

\card{SoftTint}{Card 5}
\hfill
\card{SoftBlue}{Card 6}

\newpage

\card{SoftBlue}{Card 7}
\hfill
\card{SoftTint}{Card 8}

\vspace{0.55em}

\card{SoftTint}{Card 9}
\hfill
\card{SoftBlue}{Card 10}

\section{One-Page Summary Alternative}

\noindent\colorbox{SoftTint}{%
\begin{minipage}{0.965\textwidth}
\sffamily
\textbf{\textcolor{HeadingBlue}{Rule for the one-page summary.}}
Do not copy sentences from the study pack. Use memory first, then correct with
another color. Each box needs terms, one named example or case-study hook, and
one evaluative point.
\end{minipage}
}

\vspace{0.6em}

\begin{tabularx}{\textwidth}{Y Y}
\toprule
\textbf{\textcolor{HeadingBlue}{Urban Pattern}} &
\textbf{\textcolor{HeadingBlue}{Urban Change}} \\
\midrule
\rule{0pt}{5.2cm} & \rule{0pt}{5.2cm} \\
\bottomrule
\end{tabularx}

\vspace{0.7em}

\begin{tabularx}{\textwidth}{Y Y}
\toprule
\textbf{\textcolor{HeadingBlue}{Population Process}} &
\textbf{\textcolor{HeadingBlue}{Population Policy}} \\
\midrule
\rule{0pt}{5.2cm} & \rule{0pt}{5.2cm} \\
\bottomrule
\end{tabularx}

\newpage

\section{Population Visual-Stimulus Drill}

Use this page for the Paper 2 habit: move from \textbf{what the data shows} to
\textbf{why it matters}. Do the sketch and questions from memory first, then
correct with the study pack.

\subsection{Blank Population Pyramid Template}

\begin{center}
\fcolorbox{RuleGray}{white}{%
\begin{minipage}{0.92\textwidth}
\centering
\setlength{\unitlength}{1mm}
\begin{picture}(150,58)
  \put(75,7){\line(0,1){44}}
  \put(15,7){\line(1,0){120}}
  \put(70,2){\sffamily\small 0}
  \put(18,2){\sffamily\small Male}
  \put(123,2){\sffamily\small Female}
  \put(78,47){\sffamily\small 80+}
  \put(78,37){\sffamily\small 60--79}
  \put(78,27){\sffamily\small 40--59}
  \put(78,17){\sffamily\small 20--39}
  \put(78,8){\sffamily\small 0--19}
  \put(16,52){\sffamily\small Sketch a youthful or ageing structure here}
\end{picture}
\end{minipage}
}
\end{center}

\subsection{Mini Data Task}

\rowcolors{2}{TableStripe}{white}
\begin{tabularx}{\textwidth}{L{0.22\textwidth}Y Y}
\toprule
\textbf{Age group} & \textbf{Population A} & \textbf{Population B} \\
\midrule
0--14 & 36\% & 12\% \\
15--64 & 59\% & 59\% \\
65+ & 5\% & 29\% \\
\bottomrule
\end{tabularx}
\rowcolors{2}{}{}

\begin{enumerate}
  \item Describe one clear difference between Population A and Population B.\\[0.45em]
  \rule{0.88\linewidth}{0.4pt}
  \item Explain one likely cause of Population B's structure.\\[0.45em]
  \rule{0.88\linewidth}{0.4pt}
  \item Suggest one consequence for government planning in Population A.\\[0.45em]
  \rule{0.88\linewidth}{0.4pt}
  \item Name one policy response and one limitation.\\[0.45em]
  \rule{0.88\linewidth}{0.4pt}
\end{enumerate}

\newpage

\section{Self-Check}

Score each line quickly after completing the flashcards or summary.

\rowcolors{2}{TableStripe}{white}
\begin{tabularx}{\textwidth}{Y L{0.16\textwidth} L{0.31\textwidth}}
\toprule
\textbf{Check} & \textbf{Score} & \textbf{Fix needed} \\
\midrule
I can name the four Urban environments clusters. & \_\_ /5 & \\
I can explain urbanization, suburbanization, and gentrification without notes. & \_\_ /5 & \\
I can separate urban environmental stress from social stress. & \_\_ /5 & \\
I can explain population distribution using physical and human factors. & \_\_ /5 & \\
I can use births, deaths, and migration to explain population change. & \_\_ /5 & \\
I can read a population pyramid using shape, process, consequence, and policy. & \_\_ /5 & \\
I can explain dependency ratio and link it to service or tax pressure. & \_\_ /5 & \\
I can explain demographic dividend and the conditions needed for it. & \_\_ /5 & \\
I can evaluate one population policy. & \_\_ /5 & \\
\bottomrule
\end{tabularx}
\rowcolors{2}{}{}

\section{Error Log}

Write one mistake from today's work, the corrected idea, and the next action.

\rowcolors{2}{TableStripe}{white}
\begin{tabularx}{\textwidth}{Y Y Y}
\toprule
\textbf{Mistake or weak point} & \textbf{Correction} & \textbf{Next action} \\
\midrule
\rule{0pt}{1.2cm} & & \\
\rule{0pt}{1.2cm} & & \\
\bottomrule
\end{tabularx}
\rowcolors{2}{}{}

\section{Carryover To April 25}

April 25 is Urban environments depth. Before ending today, update the 1--5
mastery ratings from the April 23 board, then choose the three weakest urban
items to drill tomorrow.

\begin{itemize}
  \item Urban environments rating after today: \rule{0.12\textwidth}{0.4pt} /5
  \item Population rating after today: \rule{0.12\textwidth}{0.4pt} /5
\end{itemize}

\begin{enumerate}
  \item Weakest urban item: \rule{0.72\textwidth}{0.4pt}
  \item Second weakest urban item: \rule{0.66\textwidth}{0.4pt}
  \item Third weakest urban item: \rule{0.68\textwidth}{0.4pt}
\end{enumerate}

\vspace{0.5em}

\noindent Main reason for weakness:
\quad definition \quad / \quad process \quad / \quad example \quad / \quad data
\quad / \quad evaluation

\end{document}
